\section{Conditioning of companion eigenvalues}
\label{sec:companion}

We now keep the scalar output fixed and change its representation.  Let
\[
  p(x)=\sum_{j=0}^{n}a_jx^j,
  \qquad a_0,\ldots,a_n\ \text{independent with law }\CN(0,1),
\]
and let $C_p$ be multiplication by $x$ on $\C[x]/(p)$ in the ordered
basis $1,x,\ldots,x^{n-1}$.  Thus $C_p$ has ones on its subdiagonal and
last column $-(a_0,\ldots,a_{n-1})^{T}/a_n$.  Almost surely $a_n\ne0$
and all roots are simple.  For a root $z$, put
\[
  v(z)=(1,z,\ldots,z^{n-1})^T,
  \qquad
  \ell_z(x)=\frac{p(x)}{(x-z)p'(z)}.
\]
The coefficient vector of $\ell_z$ is a right eigenvector of $C_p$,
whereas $v(z)^T$ is a left eigenvector, and $\ell_z(z)=1$.  Consequently
the absolute eigenvalue condition number for spectral-norm matrix
perturbations is
\begin{equation}\label{eq:comp-overlap}
  \Lambda(p,z)=\|v(z)\|_2\,\|\ell_z\|_2.
\end{equation}
Here the second norm is the Euclidean norm of the coefficient vector in
the displayed monomial basis.  In particular $\Lambda\ge1$.  This basis
choice is part of the theorem: the overlap is not invariant under a
general change of polynomial basis.

The next theorem proves part (i) of Theorem~\ref{thm:intro-main}.
The comparison in part (ii) follows in Corollary~\ref{cor:comp-relative}.

\begin{theorem}[Uniform companion-overlap tails]\label{thm:comp-main}
There is an absolute constant $C$ such that, for every $n\ge1$ and $t>0$,
\begin{equation}\label{eq:comp-weak-two}
  \Pstar\{\Lambda(p,Z)>t\}
  =\frac1n\E\#\{z:p(z)=0,\ \Lambda(p,z)>t\}
  \le \min\{1,Ct^{-2}\}.
\end{equation}
For every $0<q<2$ there is a constant $C_q$, independent of $n$, with
\[
  \Estar\Lambda^q\le C_q,
  \qquad 0\le\Estar\log\Lambda\le C.
\]
The moment range is sharp for bounds uniform in the degree:
\begin{equation}\label{eq:comp-second-asymptotic}
  \Estar\Lambda^2=2\log n+O(1)
  \qquad(n\to\infty).
\end{equation}
For $n=1$, the overlap is identically one.
\end{theorem}

The tail controls the expected fraction of poorly conditioned
eigenvalues.  It does not require independence of the roots and does not
assert a degree-independent bound for the worst root of a polynomial.
The second-moment assertion also shows that the exponent $2$ in
\eqref{eq:comp-weak-two} cannot be replaced by any larger exponent with
a degree-independent constant.

\subsection{An exact conditional moment}

Write $k_j(r)=\sum_{i=0}^{j}r^i$ and
$h_n(r)=(r k_n'(r)/k_n(r))'$, as in the box model.  At a fixed location
$z$, coefficient multiplication by $x-z$ is the injective map
\[
  D_z:\C^n\longrightarrow\C^{n+1},
  \qquad \mathbf q\longmapsto\operatorname{coeff}((x-z)q(x)).
\]
Its Gram matrix $T_z=D_z^*D_z$ is tridiagonal, with diagonal $1+r$,
superdiagonal $-z$, and subdiagonal $-\overline z$, where $r=|z|^2$.
In particular it is positive definite, including at $z=0$.

\begin{proposition}[Conditional overlap identity]\label{prop:comp-exact}
Under the selected-root law, the conditional second moment is
\begin{equation}\label{eq:comp-exact}
  \Phi_n(r):=\Estar[\Lambda^2\mid Z=z]
  =\frac{k_{n-1}(r)}{k_n(r)h_n(r)}\tr(T_z^{-1})
  =\frac{k_{n-1}(r)A_n(r)}{k_n(r)^2h_n(r)},
\end{equation}
where
\begin{equation}\label{eq:comp-trace-poly}
  A_n(r)=\sum_{j=0}^{n-1}(j+1)(n-j)r^j.
\end{equation}
The identity defines a continuous version of the conditional moment.
It satisfies
\[
  \Phi_n(r)=\Phi_n(r^{-1})\quad(r>0),
  \qquad \Phi_n(0)=n,
  \qquad \Phi_n(1)=\frac{2n}{n+1}.
\]
For $n=1$, $\Phi_1(r)=1$ for all $r\ge0$.
\end{proposition}

\begin{proof}
Ordinary Gaussian conditioning on $p(z)=0$ restricts the standard
Gaussian coefficient vector to the hyperplane $\operatorname{im}D_z$.
Writing $p(x)=(x-z)q(x)$ therefore gives
\[
  \mathbf q\mid\{p(z)=0\}\sim\CN(0,T_z^{-1}).
\]
Indeed its density in these coordinates is proportional to
$\exp(-\mathbf q^*T_z\mathbf q)$; the Jacobian of the linear
parametrization is constant.  Gaussian regression also gives
\[
  \E[|p'(z)|^2\mid p(z)=0]=k_n(r)h_n(r).
\]
At a selected root, the conditional law is tilted by $|p'(z)|^2$.
Since $p'(z)=q(z)$ and
$\Lambda^2=k_{n-1}(r)\|\mathbf q\|_2^2/|p'(z)|^2$, the tilt cancels
the denominator.  Thus
\[
  \Estar[\Lambda^2\mid Z=z]
  =\frac{k_{n-1}(r)\E[\|\mathbf q\|_2^2\mid p(z)=0]}
  {\E[|p'(z)|^2\mid p(z)=0]},
\]
which proves the first equality in \eqref{eq:comp-exact}.

The determinant of the $j\times j$ leading block of $T_z$ is $k_j(r)$:
both satisfy the recurrence
$d_j=(1+r)d_{j-1}-r d_{j-2}$, with $d_0=1$ and $d_1=1+r$.
Deleting its $i$th row and column leaves two such blocks, so
\begin{equation}\label{eq:comp-trace-minors}
  \tr(T_z^{-1})
  =\frac1{k_n(r)}\sum_{i=1}^{n}k_{i-1}(r)k_{n-i}(r)
  =\frac{A_n(r)}{k_n(r)}.
\end{equation}
The coefficient of $r^j$ in the numerator is $(j+1)(n-j)$,
which proves the second equality.

Reversal provides a useful check and removes the distinction between
inner and outer roots.  For $z\ne0$, let $p^{\mathrm{rev}}(x)=x^np(1/x)$.
Its coefficients have the same law as those of $p$, and direct reversal
of $\ell_z$ gives
\[
  \|\ell^{\mathrm{rev}}_{1/z}\|_2
   =|z|^{n-1}\|\ell_z\|_2,
  \qquad
  \|v(1/z)\|_2=|z|^{-(n-1)}\|v(z)\|_2.
\]
Hence $\Lambda(p^{\mathrm{rev}},1/z)=\Lambda(p,z)$.  Reversal carries a
uniformly selected root to a uniformly selected root, proving the
symmetry of $\Phi_n$.  Finally,
\[
  h_n(1)=\frac{n(n+2)}{12},
  \qquad \tr(T_1^{-1})=\frac{n(n+2)}6.
\]
Substitution proves the values at $r=1$; the assertions at $r=0$ and
$n=1$ follow directly from \eqref{eq:comp-exact}.
\end{proof}

\subsection{Uniform estimates and a tail bound}

For $0\le r\le1$, the cofactor formula
\eqref{eq:comp-trace-minors} immediately gives
\begin{equation}\label{eq:comp-trace-bound}
  \tr(T_z^{-1})
  \le n\min\{n,(1-r)^{-1}\}
  \le\frac{2n}{1-r+1/n},
\end{equation}
with the usual interpretation at $r=1$.  Combining this with
the metric estimate~\eqref{eq:box-metric-bound}, $k_{n-1}\le k_n$, and reversal proves
\begin{equation}\label{eq:comp-depth-bound}
  \Phi_n(r)\le C\bigl(1+n(1-\min\{r,r^{-1}\})\bigr)
  \le C(1+|n\log r|),\qquad r>0.
\end{equation}
The conditional second moment alone gives all the uniform moments of
order below two.  A little more of the same Gaussian regression gives
the sharper tail in Theorem~\ref{thm:comp-main}.

\begin{lemma}[Conditional moments]\label{lem:comp-third}
For each $0<q<4$ there is a constant $C_q$, independent of $n$ and $z$, such that
\[
  \Estar[\Lambda^q\mid Z=z]\le C_q\Phi_n(|z|^2)^{q/2}.
\]
\end{lemma}

\begin{proof}
Under ordinary conditioning on $p(z)=0$, put $\Sigma=T_z^{-1}$ and
$\sigma^2=v(z)^T\Sigma\overline{v(z)}=k_n(r)h_n(r)$.  Gaussian
regression decomposes the coefficient vector as
\[
  \mathbf q=b\xi+\eta,
  \qquad b=\Sigma\overline{v(z)}/\sigma,
  \qquad \xi=q(z)/\sigma\sim\CN(0,1),
\]
where $\eta$ is independent of $\xi$ and has covariance
$R=\Sigma-bb^*$.  At a selected root, the tilt is $|\xi|^2$.
It preserves the law of $\eta$ and its independence from $\xi$, while
$|\xi|^2$ now has density $u e^{-u}$ on $(0,\infty)$.  In particular
\[
  \Estar|\xi|^{-q}
    =\int_0^\infty u^{1-q/2}e^{-u}\,du
    =\Gamma(2-q/2)<\infty,\qquad 0<q<4.
\]
Writing $\beta=\|v(z)\|_2/\sigma$, we have
\[
  \Lambda=\beta\frac{\|b\xi+\eta\|_2}{|\xi|}
  \le\beta\|b\|_2+\beta\frac{\|\eta\|_2}{|\xi|},
  \qquad \Phi_n(r)=\beta^2\tr\Sigma.
\]
The standard Gaussian fourth-moment identity, valid also for a
degenerate covariance, is
\[
  \E\|\eta\|_2^4=(\tr R)^2+\tr(R^2)\le2(\tr R)^2.
\]
It follows by diagonalizing $R$ and using $\E|g|^2=1$ and
$\E|g|^4=2$ for a standard complex Gaussian $g$.
Hence $\E\|\eta\|_2^q\le2^{q/4}(\tr R)^{q/2}$ for $0<q<4$.  Apply
$(a+b)^q\le2^{\max\{q-1,0\}}(a^q+b^q)$, independence, and
$\|b\|_2^2\le\tr\Sigma$, $\tr R\le\tr\Sigma$ to obtain the claim.
\end{proof}

\begin{proposition}[Fixed-degree integrability]\label{prop:comp-fixed-degree}
For each fixed $n\ge2$ and $q>0$,
\[
  \Estar\Lambda^q<\infty\quad\Longleftrightarrow\quad q<4.
\]
For $n=1$ every moment is finite, since $\Lambda=1$.
\end{proposition}

\begin{proof}
The stronger first bound in \eqref{eq:comp-depth-bound} gives
$\Phi_n(r)\le Cn$ at every radius.  Lemma~\ref{lem:comp-third} therefore
proves finiteness for $q<4$.  For the converse, the residual covariance
$R$ in its proof has rank $n-1$: regression removes one complex
direction from the positive definite covariance $\Sigma$.  If $n\ge2$,
there is thus an $a>0$ with $\Prob\{\|\eta\|_2>a\}>0$.  On this event,
for sufficiently small $|\xi|$, one has $\|b\xi+\eta\|_2>a/2$.
Independence and the density $u e^{-u}$ of $|\xi|^2$ then show that
$\Estar[\Lambda^q\mid Z=z]=\infty$ for $q\ge4$, because
$\int_0^\varepsilon u^{1-q/2}e^{-u}\,du=\infty$.  Averaging over $Z$
proves the assertion.
\end{proof}

Thus the threshold four for each fixed degree comes from the small
derivative direction in the root-conditioned Gaussian.  The lower
threshold two for bounds uniform in degree also reflects the
distribution of root locations.

\begin{proof}[Proof of the tail and moment assertions in
Theorem~\ref{thm:comp-main}]
Let $D=|n\log|Z|^2|$ and $W=1+D$.  The box root law~\eqref{eq:box-depth} gives
$\Pstar\{D>u\}\le\min\{1,2/u\}$, so
$\Pstar\{W>u\}\le\min\{1,4/u\}$.  By conditional Markov,
Lemma~\ref{lem:comp-third} with $q=3$, and \eqref{eq:comp-depth-bound},
\[
  \Pstar\{\Lambda>t\}
  \le\Estar\min\{1,CW^{3/2}/t^3\}.
\]
For every $A>0$, integration of the tail gives
\[
  \Estar\min\{1,(W/A)^{3/2}\}
  =\frac{3}{2A^{3/2}}\int_0^A
       u^{1/2}\Pstar\{W>u\}\,du
  \le\min\{1,C/A\}.
\]
Taking $A$ to be a constant multiple of $t^2$ proves
\eqref{eq:comp-weak-two}.  Tail integration then proves
$\sup_n\Estar\Lambda^q<\infty$ for $0<q<2$ and
$\sup_n\Estar\log\Lambda<\infty$.  Nonnegativity of the logarithm
follows from $\Lambda\ge1$.
\end{proof}

\subsection{The critical second moment}

\begin{proposition}\label{prop:comp-second}
For every $n\ge1$,
\begin{equation}\label{eq:comp-second-integral}
  \Estar\Lambda^2
  =\frac2n\int_0^1
       \frac{k_{n-1}(r)A_n(r)}{k_n(r)^2}\,dr.
\end{equation}
In particular $\Estar\Lambda^2=2\log n+O(1)$ as $n\to\infty$.
\end{proposition}

\begin{proof}
The density of $|Z|^2$ is $h_n(r)/n$.  It cancels the $h_n$ in
\eqref{eq:comp-exact}, and reversal makes the contributions of
$(0,1)$ and $(1,\infty)$ equal.  This proves
\eqref{eq:comp-second-integral}, including finiteness for each $n$.

For $0\le r<1$, summing \eqref{eq:comp-trace-poly} gives
\[
  A_n(r)=
  \frac{n-(n+2)r+(n+2)r^{n+1}-nr^{n+2}}{(1-r)^3}.
\]
If $G_n(r)=k_{n-1}(r)A_n(r)/(n k_n(r)^2)$, rearrangement yields
\begin{align*}
  G_n(r)
  &=\frac1{1-r}
    \frac{(1-r^n)(1+r^{n+1})}{(1-r^{n+1})^2}\\
  &\hspace{1em}
    -\frac{2r}{n(1-r)^2}
      \left(\frac{1-r^n}{1-r^{n+1}}\right)^2.
\end{align*}
Suppose $n\ge2$.  On $0\le r\le1-1/n$ we have
$r^{n+1}\le e^{-1}$, and therefore
\[
  \left|G_n(r)-\frac1{1-r}\right|
  \le\frac{Cr^n}{1-r}+\frac{2}{n(1-r)^2}.
\]
The right side has a uniformly bounded integral on this interval:
for its first term use $r^n\le e^{-n(1-r)}$ and
$u=n(1-r)\ge1$; the second term integrates directly.  Hence
\[
  \int_0^{1-1/n}G_n(r)\,dr=\log n+O(1).
\]
On the remaining interval, \eqref{eq:comp-trace-bound} and
$k_{n-1}\le k_n$ give $0\le G_n(r)\le n$, so its integral is at most
one.  Multiplication by two proves the asserted asymptotic.
\end{proof}

This completes Theorem~\ref{thm:comp-main}.  By monotonicity of moments
on a probability space, $\sup_n\Estar\Lambda^q=\infty$ for every
$q\ge2$, allowing the value $+\infty$ for a fixed degree.  Thus the
sharpness assertion concerns uniformity in $n$; it does not say that
all moments of order at least two diverge at each fixed degree.

\begin{example}[Degrees one and two]\label{ex:comp-small}
When $n=1$, $C_p$ is a scalar and $\Lambda=1$.  When $n=2$,
\[
  T_z=\begin{pmatrix}1+r&-z\\-\overline z&1+r\end{pmatrix},
  \qquad
  \Phi_2(r)=\frac{2(1+r)^2}{1+4r+r^2}.
\]
In particular the conditional second moment is $4/3$ on the unit
circle, and direct integration of \eqref{eq:comp-second-integral}
gives
\[
  \Estar\Lambda^2=\frac23+\frac{4\pi}{9\sqrt3}.
\]
The deterministic overlap can nevertheless become arbitrarily large:
if $p(x)=(x-z)(x-w)$ with $z\ne w$, then
$\Lambda(p,z)=\sqrt{(1+|z|^2)(1+|w|^2)}/|z-w|$.
The theorem describes how averaging over the polynomial and a selected
root controls these near-collision configurations.
\end{example}

\subsection{Comparison with coefficient conditioning}

For the same scalar output $z$, coefficient and matrix perturbations
give, respectively,
\[
  \kappa_{\mathrm{coeff}}(p,z)
     =\frac{\sqrt{k_n(|z|^2)}}{|p'(z)|},
  \qquad
  \kappa_{\mathrm{eig}}(C_p,z)=\Lambda(p,z).
\]
All $n+1$ polynomial coefficients are perturbable in the first
condition number, including the leading coefficient.  Multiplication
by the corresponding data norm and division by the common output norm
$|z|$ define their relative versions; $z\ne0$ almost surely.

\begin{corollary}[A relative conditioning comparison]\label{cor:comp-relative}
With Euclidean coefficient norm and spectral matrix norm, define
\[
  \kappa_{\mathrm{coeff,rel}}(p,z)
    =\frac{\sqrt{k_n(|z|^2)}\|\mathbf a\|_2}{|p'(z)|\,|z|},
  \qquad
  \kappa_{\mathrm{eig,rel}}(C_p,z)
    =\frac{\Lambda(p,z)\|C_p\|_2}{|z|}.
\]
Then
\begin{align*}
  \Estar\log\frac{\kappa_{\mathrm{eig}}}{\kappa_{\mathrm{coeff}}}
    &=\log n+O(1),\\
  \Estar\log\frac{\kappa_{\mathrm{eig,rel}}}
                          {\kappa_{\mathrm{coeff,rel}}}
    &=\log n+O(1).
\end{align*}
The error terms are bounded uniformly in the degree.
\end{corollary}

\begin{proof}
The normalized derivative
$p'(Z)/\sqrt{k_n(|Z|^2)h_n(|Z|^2)}$ has the determinant-weighted
complex Gaussian law from the local theorem: its squared modulus has
density $u e^{-u}$.  Consequently the following finite-degree identity
holds, where $\gamma$ is Euler's constant:
\begin{equation}\label{eq:comp-absolute-identity}
  \Estar\log\frac{\kappa_{\mathrm{eig}}}{\kappa_{\mathrm{coeff}}}
  =\Estar\log\Lambda
    +\frac12\Estar\log h_n(|Z|^2)+\frac{1-\gamma}{2}.
\end{equation}
The bounded logarithmic overlap and the box-model estimate~\eqref{eq:box-metric-log}
$\Estar\log h_n(|Z|^2)=2\log n+O(1)$ prove the first assertion.

For the relative comparison the output factors cancel exactly.  It
remains to control
$\E\log\|C_p\|_2-\E\log\|\mathbf a\|_2$; a statistic depending only
on the polynomial has the same expectation under $\Pstar$, because
there are always $n$ roots.  With
$b=(a_0,\ldots,a_{n-1})^T/a_n$, the last column gives the lower bound
below, while $C_pC_p^*=\diag(0,1,\ldots,1)+bb^*$ gives the upper bound:
\[
  \|b\|_2\le\|C_p\|_2\le\sqrt{1+\|b\|_2^2}.
\]
Since sums of $m$ independent squared standard complex Gaussians have
the $\operatorname{Gamma}(m,1)$ law, their logarithmic means give
\begin{equation}\label{eq:comp-data-scale}
  \frac\gamma2-\frac1{2n}
  \le \E\log\|C_p\|_2-\E\log\|\mathbf a\|_2
  \le\frac\gamma2.
\end{equation}
For example, this follows from
$\E\log\operatorname{Gamma}(m,1)=H_{m-1}-\gamma$, with $H_0=0$.
Adding \eqref{eq:comp-data-scale} proves the relative assertion.
\end{proof}

These are signed logarithmic comparisons of two specified local
condition numbers.  They say that the geometric mean of their ratio is
bounded above and below by constant multiples of $n$.  The comparison
does not include errors in forming $C_p$ or in recovering information
from an eigenvector; an end-to-end stability theorem would also have to
analyze those maps.
